% !TEX root = ../algebra_02.tex

\section{Разложение конечной циклической группы}

\begin{Theorem}{Классификация циклических групп}{thm:cyclic-classification}
	Пусть $G=\langle g\rangle$ --- циклическая группа.
	\begin{enumerate}
		\item Если $|G|=\infty$, то $G\cong \mathbb{Z}$.
		\item Если $|G|=n<\infty$, то $G\cong \mathbb{Z}/n\mathbb{Z}$.
	\end{enumerate}
\end{Theorem}

\begin{proof}
	Рассмотрим гомоморфизм
	\[
	\varphi\colon \mathbb{Z} \to G, \qquad k\mapsto g^k.
	\]
	Он сюръективен, поскольку $G=\langle g\rangle$. По теореме о подгруппах группы $\mathbb{Z}$ существует $n\ge 0$ такое, что $\Ker\varphi=n\mathbb{Z}.$
	
	Если $n=0$, то $\Ker\varphi=\{0\}$, и потому $\varphi$ инъективен. Следовательно, $G\cong \mathbb{Z}$.
	
	Если $n>0$, то по первой теореме об изоморфизме
	\[
	\mathbb{Z}/n\mathbb{Z} = \mathbb{Z}/\Ker\varphi \cong \Im\varphi = G.
	\]
	Если при этом $|G|<\infty$, то обязательно $n=|G|$.
\end{proof}

\begin{Theorem}{Разложение конечной циклической группы}{thm:cyclic-decomposition}
	Пусть $m, n \in \mathbb{N}$ и $\gcd(m, n) = 1$. Тогда имеет место изоморфизм:
	\[
	\mathbb{Z}/mn\mathbb{Z} \cong (\mathbb{Z}/m\mathbb{Z}) \times (\mathbb{Z}/n\mathbb{Z}).
	\]
\end{Theorem}

\begin{proof}
	Рассмотрим элемент $g = ([1]_m, [1]_n) \in (\mathbb{Z}/m\mathbb{Z}) \times (\mathbb{Z}/n\mathbb{Z})$. Найдем его порядок.
	
	Условие $kg = 0$ в аддитивной записи означает, что $k[1]_m = [0]_m$ и $k[1]_n = [0]_n$, что равносильно $m \mid k$ и $n \mid k$. 
	
	Так как числа взаимно просты ($\gcd(m, n) = 1$), это эквивалентно делимости на их произведение: $mn \mid k$.
	Значит, наименьшее натуральное $k$, зануляющее $g$, равно $mn$, то есть порядок элемента $g$ равен $mn$.
	
	Поскольку порядок всей группы $(\mathbb{Z}/m\mathbb{Z}) \times (\mathbb{Z}/n\mathbb{Z})$ также равен $mn$, элемент $g$ порождает всю группу:
	\[
	\langle ([1]_m, [1]_n) \rangle = (\mathbb{Z}/m\mathbb{Z}) \times (\mathbb{Z}/n\mathbb{Z}).
	\]
	Так как группа порождается одним элементом, она является циклической. Любая конечная циклическая группа порядка $mn$ изоморфна $\mathbb{Z}/mn\mathbb{Z}$. Следовательно, изоморфизм доказан.
\end{proof}